\chapter{Development process}
\section{Development setup}

\subsection*{Prerequisites}

\begin{itemize}
  \item Node.js 22.13.0 or newer (the audit used 22.19.0); \Code{node:sqlite} is required.
  \item npm and the committed \Code{package-lock.json} (lockfile version 3).
  \item A modern browser for the two web views. Internet is needed only to exercise YouTube or download dependencies/tools.
  \item Windows is required for the configured NSIS packaging target and PowerShell announcement generator; ordinary web development is not explicitly restricted by configuration.
\end{itemize}

\subsection*{Install and run both development processes}

\begin{verbatim}
npm ci
npm run dev
\end{verbatim}

This starts Nest at \Code{http://127.0.0.1:3001} and Vite at \Code{http://127.0.0.1:5173}. Open the latter root for the cashier and \Code{/pantalla} for the public display. Vite proxies API/media/Socket.IO to Nest.

Separate commands are available:
\begin{verbatim}
npm run dev:server
npm run dev:client
\end{verbatim}

\subsection*{Database and local media}

No manual initialization is required. The server creates \Code{data/comal.sqlite}, tables, indexes, and built-in music directories at startup. Override data locations with environment variables when isolation is needed. Committed demo music is already available.

To add local music, put supported files below \Code{data/music/<genre>/} and invoke refresh from the UI. Supported extensions are MP3, WAV, OGG, M4A, AAC, and FLAC.

\subsection*{Checks}

\begin{verbatim}
npm test
npx tsc --noEmit -p tsconfig.json
npx tsc --noEmit -p tsconfig.server.json
node tooling/audio/verify-audio.mjs
\end{verbatim}

There is no lint or format command. All four checks were validated during this audit.

\subsection*{Build and run}

\begin{verbatim}
npm run build
npm start
\end{verbatim}

The build checks frontend types, creates \Code{dist/}, and compiles the server into \Code{build/}. The start command serves the combined app at loopback port 3001 by default. \Code{npm run start:server} runs the TypeScript server directly and expects an existing \Code{dist/} if the SPA is needed.

For desktop development after a build:
\begin{verbatim}
npm run build
npm run desktop
\end{verbatim}

For Windows artifacts:
\begin{verbatim}
npm run desktop:dir
npm run desktop:build
\end{verbatim}

These run electron-builder; signing, installer acceptance, and clean-machine verification are not configured/documented.

\subsection*{Regenerating optional audio assets}

Runtime uses committed WAV files and does not require a synthesizer. On Windows, announcements can be regenerated with an existing eSpeak NG executable:
\begin{verbatim}
pwsh -ExecutionPolicy Bypass -File \
  tooling/audio/generate-announcements.ps1 \
  -EspeakPath 'C:\path\to\espeak-ng.exe'
\end{verbatim}

The optional \Code{-PrepareEngine} downloads pinned eSpeak NG 1.52.0, verifies its SHA-256, and extracts it locally using 7-Zip; it requires Internet but does not install the engine. Demo music is deterministic:
\begin{verbatim}
node tooling/audio/generate-demo-music.mjs
node tooling/audio/verify-audio.mjs
\end{verbatim}

Regeneration changes committed binary assets and should be reviewed deliberately.
